% ECONOMICS_PROBLEM: {"id":"C32-80","title":"Uniform inference near higher-moment identification failure","jel_code":"C32","source_id":"OP-0603"}
% FIRST_POSTED: 2026-10-10
% ENTRY_KIND: research_attempt
% ATTEMPT_STATUS: solved
\documentclass[11pt,a4paper]{article}
\usepackage[margin=25mm]{geometry}
\usepackage{amsmath,amssymb,amsthm}
\usepackage[hidelinks]{hyperref}
\hypersetup{pdftitle={Uniform Inference with Jointly Weak Variance and Fourth-Moment Features},pdfauthor={Ji Ho Bae}}
\setlength{\parindent}{0pt}
\setlength{\parskip}{4pt}
\setlength{\emergencystretch}{2em}
\allowdisplaybreaks[2]
\newtheorem{theorem}{Theorem}
\numberwithin{equation}{section}
\newcommand{\R}{\mathbb R}
\newcommand{\E}{\mathbb E}
\newcommand{\Prob}{\mathbb P}
\newcommand{\cM}{\mathcal M}
\newcommand{\cG}{\mathcal G}
\newcommand{\cX}{\mathcal X}
\newcommand{\T}{\mathsf T}
\newcommand{\F}{\mathrm F}
\newcommand{\op}{\mathrm{op}}
\newcommand{\norm}[1]{\left\lVert #1\right\rVert}
\DeclareMathOperator{\diam}{diam}
\DeclareMathOperator{\diag}{diag}
\DeclareMathOperator{\tr}{tr}
\DeclareMathOperator{\TV}{TV}
\DeclareMathOperator{\KL}{KL}
\DeclareMathOperator{\cum}{cum}
\title{Uniform Inference with Jointly Weak\\
Variance and Fourth-Moment Features}
\author{Ji Ho Bae\thanks{Prepared with GPT Codex assistance. The arguments are hand derivations. No Lean formalization or external peer review is claimed, and no novelty claim is made.}}
\date{10 October 2026}
\begin{document}
\maketitle
\begin{center}
\small
\textbf{C32-80 \quad JEL C32 \quad Source OP-0603}\\
Complete hand proof with internal mathematical and scope review.
\end{center}
\begin{abstract}
We study a declared nonparametric experiment with independent observations
from known volatility regimes, a common nonsingular impact matrix,
conditionally independent shocks, and bounded eighth moments.
Variance contrasts and baseline fourth cumulants may weaken jointly.
A global inverse inequality yields uniformly honest confidence sets for
the full impact matrix modulo signed permutations. We prove matching
worst-case expected-diameter bounds at every confidence level, including
complete feature failure and arbitrary weakening sequences. A single
procedure also attains the rate without knowing the strength lower bound.
We give matching squared-error bounds and state separately the conditions
for labelled matrices, other responses, and dimension one.
\end{abstract}

\section{Declared experiment, target, and loss}\label{sec:model}
Fix integers \(d\ge2\) and \(S\ge0\), and constants
\[
 0<b_-<b_+<\infty,\qquad
 0<v_-<1\le v_+<\infty,\qquad
 M_8\ge105,\qquad 0<\nu\le\frac1{S+1}.
\]
All constants, including \(d\) and \(S\), remain fixed as the sample size
increases. There are \(S+1\) known regimes, numbered \(0,\ldots,S\).
For each admissible sample size \(T\), their known deterministic counts
satisfy
\[
 \sum_{s=0}^S T_s=T,\qquad T_s\ge\nu T.
\]
Observations are independent across all indices. Within regime \(s\),
they have the common distribution of \(U=BE_s\).
The coordinates of \(E_s\) are mutually independent, have mean zero,
have variances \(d_{si}\in[v_-,v_+]\), and satisfy
\[
 \E|E_{si}|^8\le M_8.
\]
In regime \(0\), \(d_{0i}=1\) for every \(i\). The matrix \(B\) is common
across regimes and every singular value lies in \([b_-,b_+]\).
Apart from these conditions, the coordinate shock laws are arbitrary and
unknown; the experiment imposes no density or parametric shape assumption.
The laws and \(B\) may depend on \(T\). The observed data \(O_T\) are the
known-regime innovation vectors \(U\), not estimated residuals from an
additional, unspecified time-series model.

Write \(D_s=\diag(d_{s1},\ldots,d_{sd})\), \(c_i=\E E_{0i}^4-3\), and set
\begin{equation}\label{eq:delta}
 \delta(D,c)^2
 =\min_{i<j}\left\{\sum_{s=1}^S(d_{si}-d_{sj})^2+c_i^2+c_j^2\right\}.
\end{equation}
The sum is zero when \(S=0\). This index describes the specified moment
features, not every possible form of non-Gaussian identification.
Let \(\cM_T(k)\) be the structural models just declared with
\(\delta(D,c)\ge k\). The nonnegative \(k\) indexes the benchmark class.
The first confidence construction uses \(k\); Section~6A gives a single
procedure without it. Below, \(k\in[0,\bar k]\), for a sufficiently small
fixed \(\bar k>0\) specified in Section~\ref{sec:hard}.
This interval includes arbitrary rates of convergence to zero and the
sequence identically zero.

Let \(\cG\) be the finite group of signed permutation matrices of order \(d\).
The primary response is the entire impact matrix modulo its column signs
and labels:
\[
 \tau(m)=[B],\qquad
 d_{\cG}([B],[B'])=\min_{P\in\cG}\norm{B-B'P}_{\F}.
\]
This is a quotient metric: the action is isometric, the minimum is attained
because \(\cG\) is finite, and composition of minimizing group elements
proves the triangle inequality. Denote the compact quotient target space
by \(\cX\). Its diameter is at most
\[
 D_{\max}=2\sqrt d\,b_+.
\]

A confidence procedure is a measurable random closed subset \(C_T\) of
\(\cX\); independent auxiliary randomization is allowed. Set
\(\diam(\varnothing)=0\). For \(0<\alpha<1\), define
\[
 \mathcal R_T(k,\alpha)
 =\inf_C\ \sup_{m\in\cM_T(k)}\E_m\diam(C),
\]
where the infimum ranges over procedures satisfying
\[
 \inf_{m\in\cM_T(k)}\Prob_m\{[B_m]\in C\}\ge1-\alpha.
\]
The full space \(\cX\) is one such procedure. Define the rate
\begin{equation}\label{eq:rate}
 r(T,k)=
 \begin{cases}
  \min\{1,(\sqrt T\,k)^{-1}\},&k>0,\\
  1,&k=0.
 \end{cases}
\end{equation}

\begin{theorem}[Full quotient-matrix response]\label{thm:main}
There is \(\bar k>0\), depending only on the displayed class constants,
such that for each \(0<\alpha<1\) there are constants
\(0<c_\alpha\le C_\alpha<\infty\) with
\begin{equation}\label{eq:mainrate}
 c_\alpha r(T,k)\le\mathcal R_T(k,\alpha)\le C_\alpha r(T,k)
\end{equation}
for every admissible \(T\) and every \(k\in[0,\bar k]\).
The upper procedure has finite-sample coverage at least \(1-\alpha\);
its diameter bound holds deterministically for every data set.
It therefore holds in worst-case expectation without an
\(\alpha\)-sized remainder. At \(k=0\), the worst-case expected diameter
is bounded away from zero.

For any sequence \(k_T\in[0,\bar k]\), every procedure sequence satisfying
\[
 \liminf_{T\to\infty}\inf_{m\in\cM_T(k_T)}
 \Prob_m\{[B_m]\in C_T\}\ge1-\alpha
\]
obeys
\begin{equation}\label{eq:asymptotic}
 \liminf_{T\to\infty}
 \frac{\sup_{m\in\cM_T(k_T)}\E_m\diam(C_T)}{r(T,k_T)}
 \ge c'_\alpha>0.
\end{equation}
The finite-sample upper procedure also attains this asymptotic rate.
\end{theorem}

The finite-\(T\) infimum is deliberately over finite-\(T\) honest
procedures. Taking a finite-\(T\) infimum over merely asymptotically
honest sequences would be invalid. Labelled matrices, response maps,
covariance responses, and \(d=1\) are treated in
Sections~\ref{sec:labels}--\ref{sec:checks}. No universal rate is asserted
for arbitrary response functionals without stated regularity and nonflatness.

\section{Raw features and a uniform concentration bound}\label{sec:raw}
Use Frobenius norms on matrices and fourth-order tensors, including all
ordered tensor entries. Give a tuple the Euclidean sum-of-squares norm.
For a symmetric matrix \(\Sigma\), define its Wick tensor by
\[
 W(\Sigma)_{abcd}
 =\Sigma_{ab}\Sigma_{cd}+\Sigma_{ac}\Sigma_{bd}
   +\Sigma_{ad}\Sigma_{bc}.
\]
For \(\vartheta=(B,D_1,\ldots,D_S,c)\), let
\begin{align*}
 \Sigma_0&=BB^\T,& \Sigma_s&=BD_sB^\T\quad(1\le s\le S),\\
 K_{\mathrm{raw}}&=\sum_i c_i b_i^{\otimes4},&
 M_4&=K_{\mathrm{raw}}+W(\Sigma_0),\\
 F(\vartheta)&=(\Sigma_0,\ldots,\Sigma_S,M_4),
\end{align*}
where \(b_i\) is column \(i\) of \(B\).
Zero means and coordinate independence give the second-moment formulas.
For the fourth-moment formula, expand a product of four coordinates of
\(BE_0\). Independence and zero means annihilate every term with a coordinate
appearing once. Pair-pair terms give \(W(\Sigma_0)\); four-equal-coordinate
terms leave coefficient \(\E E_{0i}^4-3\). Thus the formula holds for
nonsymmetric shock laws and in the presence of nonzero third moments.

Let \(\widehat F\) consist of the empirical raw second moment in each
regime and the empirical raw fourth moment in regime \(0\).
These moments are not centered by a sample mean: the experiment imposes
zero population means, so the raw statistics are unbiased.

Put \(H_{\mathrm{obs},8}=b_+^8d^4M_8\). For every regime,
\begin{align*}
 \E\norm{U}^8
 &\le b_+^8\,\E\left(\sum_i E_{si}^2\right)^4\\
 &\le b_+^8d^3\sum_i\E|E_{si}|^8
 \le H_{\mathrm{obs},8}.
\end{align*}
Here \((\sum_i a_i)^4\le d^3\sum_i a_i^4\) for \(a_i\ge0\)
follows from H\"older's inequality.
Consequently \(\E\norm U^4\le\sqrt{H_{\mathrm{obs},8}}\).
Each entry \(U_aU_b\) has variance at most \(\sqrt{H_{\mathrm{obs},8}}\);
each entry \(U_aU_bU_cU_e\) has variance at most \(H_{\mathrm{obs},8}\).
Independence across observations and \(T_s\ge\nu T\) therefore imply
\begin{equation}\label{eq:concentration}
 \E_m\norm{\widehat F-F(\vartheta_m)}^2\le\frac VT,\qquad
 V=\nu^{-1}\left[(S+1)d^2\sqrt{H_{\mathrm{obs},8}}
                      +d^4H_{\mathrm{obs},8}\right].
\end{equation}
Correlations between feature entries are harmless: the expected squared
Euclidean norm is the sum of the individual second moments.
For the confidence radius
\begin{equation}\label{eq:radius}
 R_T=\frac{\sqrt{V/\alpha}}{\sqrt T},
\end{equation}
Markov's inequality gives
\begin{equation}\label{eq:coverageevent}
 \sup_m\Prob_m\{\norm{\widehat F-F(\vartheta_m)}>R_T\}\le\alpha.
\end{equation}

\section{A compact outer parameter set}\label{sec:outer}
The true cumulants satisfy
\[
 -2\le c_i\le\sqrt{M_8}-3,
\]
because \(\E E_{0i}^4\ge1\) and
\(\E E_{0i}^4\le\sqrt{\E E_{0i}^8}\).
Use the larger compact box \(|c_i|\le K\), where \(K=3+\sqrt{M_8}\).
Let \(\Theta_k\) contain the tuples \(\vartheta=(B,D,c)\) satisfying
the singular-value and diagonal-variance bounds of Section~\ref{sec:model},
this cumulant box, and \(\delta(D,c)\ge k\).
Every true parameter in \(\cM_T(k)\) belongs to \(\Theta_k\).
The set is compact and is specified by finitely many polynomial
inequalities: positive semidefiniteness of
\(B^\T B-b_-^2I\) and \(b_+^2I-B^\T B\), the diagonal boxes, and
\[
 \sum_s(d_{si}-d_{sj})^2+c_i^2+c_j^2\ge k^2
 \quad\text{for every }i<j.
\]
Define
\begin{equation}\label{eq:confidence}
 \begin{split}
 A_T(\widehat F)
   &=\{\vartheta\in\Theta_k:\norm{F(\vartheta)-\widehat F}\le R_T\},\\
 C_T&=\{[B]:(B,D,c)\in A_T(\widehat F)\}.
 \end{split}
\end{equation}
Candidates in \(\Theta_k\) need not be realizable by actual independent
shock laws. This outer enlargement is intentional: the confidence set
need only contain the true structural parameter, and all deterministic
inverse bounds below hold for every candidate. There is no hidden
unknown-distribution feasibility condition in \eqref{eq:confidence}.

By \eqref{eq:coverageevent}, the true tuple belongs to \(A_T\) with
probability at least \(1-\alpha\). Thus \(C_T\) is uniformly honest,
including at \(k=0\). It may be empty on some data sets; the specified
zero-diameter convention applies. The construction never whitens the
empirical covariance and remains defined when that covariance is singular
or ill-conditioned.

For completeness, these sets are measurable. Their graph in
\((\widehat F,\vartheta)\) is closed and the parameter coordinate lies in a
fixed compact set. Projection after additionally restricting \([B]\) to
any closed target set remains closed. Every open set in a compact metric
space is a countable union of closed balls lying inside it; consequently
the event that \(C_T\) hits an open target set is Borel.
For \(a>0\), the set of feature vectors with \(\diam(C_T)\ge a\) is also
closed: introduce two candidate tuples, impose
\(\min_{P\in\cG}\norm{B-B'P}_{\F}\ge a\), and use compact projection.
Strict events follow by countable unions. Composing with the polynomial
empirical moment map proves measurability.
The constraints are explicit finite algebraic constraints; this statistical
theorem asserts neither efficient numerical optimization nor exact
computation for arbitrary unencoded real inputs.

\section{Uniformly Lipschitz whitening}\label{sec:whitening}
For two candidate tuples, put
\(e=\norm{F(\vartheta)-F(\vartheta')}\), and let
\[
 H=\Sigma_0^{1/2},\quad H'=(\Sigma_0')^{1/2},\qquad
 L=H^{-1},\quad L'=(H')^{-1},
\]
using principal positive definite square roots.
The eigenvalues of \(H,H'\) lie in \([b_-,b_+]\).
For \(X=H-H'\),
\[
 HX+XH'=\Sigma_0-\Sigma_0'.
\]
Taking Frobenius inner products with \(X\) gives
\[
 2b_-\norm X_{\F}^2
 \le\langle X,HX+XH'\rangle
 \le\norm X_{\F}\norm{\Sigma_0-\Sigma_0'}_{\F}.
\]
Indeed, the two inner-product terms are
\(\tr(X^\T HX)\) and \(\tr(X^\T XH')\), each at least
\(b_-\norm X_{\F}^2\). Hence, also when \(X=0\),
\begin{equation}\label{eq:sqrtlip}
 \norm{H-H'}_{\F}\le\frac{\norm{\Sigma_0-\Sigma_0'}_{\F}}{2b_-},
 \qquad
 \norm{L-L'}_{\F}\le\frac{\norm{\Sigma_0-\Sigma_0'}_{\F}}{2b_-^3}.
\end{equation}
The second bound uses \(L-L'=H^{-1}(H'-H)(H')^{-1}\).

The whitened features are
\[
 A_s=L\Sigma_sL,\quad
 K_{\mathrm w}=L^{\otimes4}(M_4-W(\Sigma_0)),\qquad
 F_{\mathrm w}(\vartheta)=(A_1,\ldots,A_S,K_{\mathrm w}).
\]
Writing \(Q=LB\), we have \(QQ^\T=I\), so \(Q\) is orthogonal and
\begin{equation}\label{eq:whitened}
 A_s=QD_sQ^\T,\qquad K_{\mathrm w}=\sum_i c_iq_i^{\otimes4}.
\end{equation}
Candidate covariances obey
\(\norm{\Sigma_0}_{\F}\le\sqrt d\,b_+^2\) and
\(\norm{\Sigma_s}_{\op}\le v_+b_+^2\).
Telescoping \(L\Sigma_sL\) and applying \eqref{eq:sqrtlip} yield
\[
 \norm{A_s-A_s'}_{\F}\le C_Ae,\qquad
 C_A=b_-^{-2}+v_+b_+^2b_-^{-4}.
\]
Each of the three summands of \(W(\Sigma)\) rearranges
\(\Sigma\otimes\Sigma\), so
\begin{align*}
 \norm{W(\Sigma_0)-W(\Sigma_0')}_{\F}
 &\le3(\norm{\Sigma_0}_{\F}+\norm{\Sigma_0'}_{\F})
                \norm{\Sigma_0-\Sigma_0'}_{\F}\\
 &\le6\sqrt d\,b_+^2e.
\end{align*}
Also \(\norm{K_{\mathrm{raw}}'}_{\F}
\le\sum_i|c_i'|\norm{b_i'}^4\le dKb_+^4\).
The four-factor telescoping identity gives
\[
 \norm{L^{\otimes4}-(L')^{\otimes4}}_{\op}
 \le4b_-^{-3}\norm{L-L'}_{\op}.
\]
Consequently
\[
 \norm{K_{\mathrm w}-K_{\mathrm w}'}_{\F}\le C_Ke,\qquad
 C_K=b_-^{-4}(1+6\sqrt d\,b_+^2)+2dKb_+^4b_-^{-6}.
\]
We have proved
\begin{equation}\label{eq:whitenlip}
 \norm{F_{\mathrm w}(\vartheta)-F_{\mathrm w}(\vartheta')}
 \le L_{\mathrm w}\norm{F(\vartheta)-F(\vartheta')},\qquad
 L_{\mathrm w}=\sqrt{SC_A^2+C_K^2}.
\end{equation}
The constants depend only on the declared compact bounds.
No local approximation, eigenvector selection, or eigenvalue-gap
assumption has been used.

\section{The global joint inverse inequality}\label{sec:inverse}
We include the finite matching argument. For diagonal matrices
\(D_s=\diag(d_{s1},\ldots,d_{sd})\), tensor
\(K_0=\sum_i c_ie_i^{\otimes4}\), and \(R\in O(d)\), put
\(P_{ij}=R_{ij}^2\). The matrix \(P\) is doubly stochastic.
Let \(\mathcal E(R)^2\) be the sum of the squared distances of
\(R^\T D_sR\) from the diagonal matrix subspace and the squared distance
of \(K_0\), expressed in the \(R\) basis, from
\(\operatorname{span}\{e_j^{\otimes4}\}\).

For each covariance feature, its off-diagonal squared norm is
\begin{equation}\label{eq:covresidual}
 \sum_i d_{si}^2-\sum_j\left(\sum_i d_{si}P_{ij}\right)^2
 =\sum_{i<\ell}(d_{si}-d_{s\ell})^2\sum_jP_{ij}P_{\ell j}.
\end{equation}
To verify this, use \(\sum_jP_{ij}=1\) and apply in each column the
weighted-variance identity
\[
 \sum_iw_ix_i^2-\left(\sum_iw_ix_i\right)^2
 =\sum_{i<\ell}w_iw_\ell(x_i-x_\ell)^2,\qquad \sum_iw_i=1.
\]
The diagonal fourth-tensor coefficients in the \(R\) basis are
\[
 \beta_j=\sum_i c_iR_{ij}^4=\sum_i c_iP_{ij}^2.
\]
Orthogonal changes of basis preserve the tensor Frobenius norm, whose
square is \(\sum_ic_i^2\). Cauchy--Schwarz and
\(\sum_iP_{ij}^2\le\sum_iP_{ij}=1\) imply
\[
 \beta_j^2
 \le\left(\sum_ic_i^2P_{ij}^2\right)\left(\sum_iP_{ij}^2\right)
 \le\sum_ic_i^2P_{ij}^2.
\]
Its squared residual is therefore at least
\begin{equation}\label{eq:tensorresidual}
 \sum_ic_i^2\left(1-\sum_jP_{ij}^2\right)
 =\sum_{i<\ell}(c_i^2+c_\ell^2)\sum_jP_{ij}P_{\ell j}.
\end{equation}
No sign restriction on \(c_i\) is used. Combining
\eqref{eq:covresidual}--\eqref{eq:tensorresidual} gives
\begin{equation}\label{eq:mixingdefect}
 \mathcal E(R)^2
 \ge\delta(D,c)^2\sum_{i<\ell}\sum_jP_{ij}P_{\ell j}
 =\frac{\delta(D,c)^2}{2}\left(d-\sum_{ij}P_{ij}^2\right).
\end{equation}

Every doubly stochastic matrix is a convex combination of permutation
matrices, as the following elementary proof shows.
In the positive-entry bipartite graph, the neighbor set \(N(I)\) of any
set \(I\) of rows satisfies
\[
 |I|=\sum_{i\in I,j}P_{ij}
 \le\sum_{j\in N(I)}\sum_iP_{ij}=|N(I)|.
\]
This condition gives a perfect matching. Indeed, take a maximum matching.
An unmatched row generates an alternating search tree. Without an
augmenting path, every reached column is matched and its matched row is
reached. The reached rows then outnumber the reached columns while all
their neighbors are reached, contradicting the displayed inequality.
Thus there is a permutation \(\pi\) whose selected entries are positive.
Subtract \(\lambda\) times its permutation matrix, where \(\lambda\) is the
smallest selected entry. If \(\lambda<1\), divide the residual by
\(1-\lambda\); it is doubly stochastic and has fewer positive entries.
If \(\lambda=1\), the original matrix is that permutation matrix.
Induction on the number of positive entries proves the decomposition.

Apply the decomposition to \(P\) and take its inner product with \(P\).
Some permutation \(\pi\) satisfies
\[
 \sum_iP_{i,\pi(i)}\ge\sum_{ij}P_{ij}^2.
\]
Choose signs matching the entries \(R_{i,\pi(i)}\).
Since \(|r|\ge r^2\) when \(|r|\le1\),
\begin{align*}
 \min_{P_0\in\cG}\norm{R-P_0}_{\F}^2
 &=2d-2\max_\pi\sum_i|R_{i,\pi(i)}|\\
 &\le2\left(d-\sum_{ij}P_{ij}^2\right).
\end{align*}
Together with \eqref{eq:mixingdefect}, this proves, for \(\delta(D,c)>0\),
\begin{equation}\label{eq:inverseorthogonal}
 \min_{P_0\in\cG}\norm{R-P_0}_{\F}
 \le\frac{2\mathcal E(R)}{\delta(D,c)}.
\end{equation}

For two tuples of the form \eqref{eq:whitened}, put \(R=Q^\T Q'\).
Express both in the \(Q'\) basis and project onto the off-diagonal matrix
and fourth-tensor subspaces. The competing tuple, with its own arbitrary
diagonal \(D_s'\) and \(c_i'\), is annihilated. Orthogonal projection is a
contraction, so
\[
 \mathcal E(Q^\T Q')
 \le\norm{F_{\mathrm w}(\vartheta)-F_{\mathrm w}(\vartheta')}.
\]
Equation~\eqref{eq:inverseorthogonal} therefore yields
\begin{equation}\label{eq:inversefeatures}
 \min_{P_0\in\cG}\norm{Q-Q'P_0}_{\F}
 \le\frac{2\norm{F_{\mathrm w}(\vartheta)-F_{\mathrm w}(\vartheta')}}
              {\delta(D,c)}.
\end{equation}
The two signed-permutation minimizations agree by orthogonal
multiplication and replacing a group element by its inverse.
This is a global inequality allowing different nuisance variances and
cumulants, not merely a derivative-rank calculation.

\section{Deterministic confidence diameter and the upper theorem}\label{sec:upper}
Reconstruct \(B=HQ\) and \(B'=H'Q'\).
Equations~\eqref{eq:sqrtlip}, \eqref{eq:whitenlip}, and
\eqref{eq:inversefeatures} imply
\begin{align}
 d_{\cG}([B],[B'])
 &\le\norm{H-H'}_{\F}
       +b_+\min_{P_0\in\cG}\norm{Q-Q'P_0}_{\F}\notag\\
 &\le\left[\frac1{2b_-}+\frac{2b_+L_{\mathrm w}}{\delta(D,c)}\right]
             \norm{F(\vartheta)-F(\vartheta')}.\label{eq:inverseB}
\end{align}
Restrict \(\bar k\le1\). For any \(\vartheta,\vartheta'\in\Theta_k\)
with \(k>0\), this gives
\begin{equation}\label{eq:inversek}
 d_{\cG}([B],[B'])
 \le\frac{L_I}{k}\norm{F(\vartheta)-F(\vartheta')},\qquad
 L_I=\frac1{2b_-}+2b_+L_{\mathrm w}.
\end{equation}
Any two accepted candidates in \(A_T\) have raw-feature distance at most
\(2R_T\). Hence, on every data set, not merely on a coverage event,
\begin{equation}\label{eq:detdiameter}
 \diam(C_T)
 \le\min\left\{D_{\max},
       \frac{2L_I\sqrt{V/\alpha}}{\sqrt T\,k}\right\}
 \quad(k>0).
\end{equation}
At \(k=0\), \(\diam(C_T)\le D_{\max}\).
Empty sets satisfy the same bounds.
For \(x\ge0\),
\(\min\{D_{\max},Ax\}\le\max\{D_{\max},A\}\min\{1,x\}\).
Together with finite-sample coverage from Section~\ref{sec:outer},
\eqref{eq:detdiameter} proves the upper bound in
Theorem~\ref{thm:main}, with
\[
 C_\alpha=\max\{D_{\max},\,2L_I\sqrt{V/\alpha}\}.
\]

Identification from these raw moments is an immediate corollary:
if two candidates have identical raw moments and the first has
\(\delta>0\), then \eqref{eq:inverseB} forces identical quotient impact
matrices. This does not assert that all distributional identification
fails whenever \(\delta=0\).

\section*{6A.\quad One confidence procedure without the strength bound}
Define \eqref{eq:confidence} using \(\Theta_0\) for every sample, without
inputting \(k\), and call its quotient image \(C_T^{\mathrm{adapt}}\).
The same concentration event gives coverage at least \(1-\alpha\)
uniformly over the entire declared structural class, including
\(\delta=0\).

Fix any true tuple \(\vartheta_*\) with \(\delta_*>0\), and put
\[
 e_*=\norm{\widehat F-F(\vartheta_*)}.
\]
Every accepted candidate \(\vartheta\) satisfies
\[
 \norm{F(\vartheta)-F(\vartheta_*)}\le R_T+e_*.
\]
The bound \eqref{eq:inverseB} uses only the separation of its first tuple.
Apply it with the true tuple first, even if \(\vartheta_*\) is not
accepted by the confidence set. Comparing each of two accepted
candidates to the true target and using the triangle inequality gives,
for every data set, with \(L_0=(2b_-)^{-1}\) and
\(L_1=2b_+L_{\mathrm w}\),
\begin{equation}\label{eq:adaptdata}
 \diam(C_T^{\mathrm{adapt}})
 \le\min\left\{D_{\max},
       2\left[L_0+\frac{L_1}{\delta_*}\right](R_T+e_*)\right\}.
\end{equation}
The candidates themselves need not have positive \(\delta\);
empty sets again satisfy the bound.
By \eqref{eq:concentration}, \(\E e_*\le\sqrt{V/T}\).
Taking expectations of each of the two bounds in
\eqref{eq:adaptdata} yields
\begin{equation}\label{eq:adaptexpectation}
 \E\diam(C_T^{\mathrm{adapt}})
 \le\min\left\{D_{\max},
 \frac{2[L_0+L_1/\delta_*]\,[\sqrt{V/\alpha}+\sqrt V]}{\sqrt T}
 \right\}.
\end{equation}
For \(\delta_*\ge k>0\) and \(k\le\bar k\le1\), this is at most a fixed
multiple of \(r(T,k)\); for \(k=0\), use \(D_{\max}\).
Thus a single procedure, without knowledge of \(k\) or \(\delta_*\),
attains the upper minimax expected-diameter rate on every \(\cM_T(k)\).
Its every-data-set bound is \eqref{eq:adaptdata}, which contains the
realized moment error. The fixed numerical bound
\eqref{eq:detdiameter} belongs to the \(k\)-restricted procedure.

\section{A non-Gaussian hard family inside the nonparametric class}
\label{sec:hard}
Let \(\varphi\) be the \(N(0,1)\) density and define
\begin{align*}
 H_4(z)&=z^4-6z^2+3,\\
 H_8(z)&=z^8-28z^6+210z^4-420z^2+105,\\
 h(x)&=e^{-x^2}\left[H_4(\sqrt3\,x)-\frac14H_8(\sqrt3\,x)\right].
\end{align*}
The functions \(h,h'\) are bounded because they are polynomials times
\(e^{-x^2}\). Put \(M_h=\norm h_\infty>0\).
For \(r=0,2,4,8\), the change of variable \(y=\sqrt3\,x\) gives
\[
 \E_\varphi[X^rh(X)]
 =3^{-(r+1)/2}\E\left[Y^r\left(H_4(Y)-\frac14H_8(Y)\right)\right],
 \qquad Y\sim N(0,1).
\]
For \(j=4,8\), direct differentiation gives
\(H_j(y)\varphi(y)=(-1)^j\varphi^{(j)}(y)\).
Repeated integration by parts, with vanishing boundary terms because a
Gaussian dominates each polynomial, therefore yields
\[
 \E[Y^rH_j(Y)]
 =\E\left[\left.\frac{d^j}{dy^j}y^r\right|_{y=Y}\right].
\]
Consequently
\begin{equation}\label{eq:hermitemoments}
 \begin{aligned}
 \E_\varphi h&=0,& \E_\varphi X^2h&=0,\\
 A:=\E_\varphi X^4h&=\frac{24}{9\sqrt3}
                    =\frac8{3\sqrt3}>0,\\
 \E_\varphi X^8h
 &=\frac{5040-40320/4}{81\sqrt3}
   =-\frac{5040}{81\sqrt3}.
 \end{aligned}
\end{equation}
For example,
\(\E[Y^8H_4(Y)]=(8!/4!)\E Y^4=5040\) and
\(\E[Y^8H_8(Y)]=8!=40320\), verifying the numerical constants.

For \(0\le\varepsilon\le(2M_h)^{-1}\), the function
\[
 f_\varepsilon(x)=\varphi(x)(1+\varepsilon h(x))
\]
is a strictly positive symmetric density between
\(\varphi/2\) and \(3\varphi/2\).
Equation~\eqref{eq:hermitemoments} shows that it has mean zero, variance
one, fourth cumulant \(A\varepsilon\), and eighth moment
\(105-5040\varepsilon/(81\sqrt3)\).
In particular,
\begin{equation}\label{eq:perturbmoments}
 \cum_4(f_{k/A})=k,\qquad
 \E_{f_{k/A}}X^8=105-\frac{70}{3}k\le105.
\end{equation}
For \(k>0\), the nonzero fourth cumulant proves non-Gaussianity.

Choose henceforth
\begin{equation}\label{eq:kbar}
 \bar k=\min\left\{1,\frac1{10d},\frac{1-v_-}{2d},
                              \frac A{2M_h}\right\}>0.
\end{equation}
For \(k\in[0,\bar k]\), let \(Z_1,\ldots,Z_d\) be independent with
density \(f_{k/A}\). In the baseline regime set \(E_0=Z\).
If \(S\ge1\), use
\[
 D_1=\diag(1-k,1-2k,\ldots,1-dk),\qquad E_1=D_1^{1/2}Z,
\]
and \(D_s=I,E_s=Z\) for \(s\ge2\), with independent copies for all
observations. When \(S=0\), there are no other regimes.
Every variance lies both in \([0.9,1]\) and in \([v_-,v_+]\).
Every coordinate eighth moment is at most \(105\), because the scale
factors are at most one. This family therefore meets the same moment
cap \(M_8\ge105\) as the upper theorem.
Its baseline cumulants are all \(k\), and
\[
 \delta(D,c)=
 \begin{cases}
  \sqrt2\,k,&S=0,\\
  \sqrt3\,k,&S\ge1.
 \end{cases}
\]
For the second case, adjacent coordinates minimize
\((i-j)^2k^2+2k^2\).
Thus the family lies in \(\cM_T(k)\), including at \(k=0\).
When \(S\ge1\), both nontrivial variance contrasts and baseline fourth
cumulants vanish with \(k\).

\section{Quantitative closeness of the rotation experiments}\label{sec:closeness}
Put \(b_0=(b_-+b_+)/2\). Let \(R_\theta\) rotate the first two coordinates
through angle \(\theta\) and act as the identity on the others.
Set \(B_\theta=b_0R_\theta\), initially for
\(|\theta|\le\theta_0=1/8\). Every singular value equals \(b_0\), so the
general spectral bounds are respected; they need not contain \(1\).
Let \(P_{\theta,s}\) be the one-observation law in regime \(s\) with the
shocks of Section~\ref{sec:hard}.
A common rescaling \(x=b_0y\) leaves chi-square divergence unchanged.
In the \(y\) coordinates, the density is
\begin{equation}\label{eq:rotationdensity}
 \begin{aligned}
 p_{\theta,s}(y)&=g_{\theta,s}(y)H_{\theta,s}(y),\\
 g_{\theta,s}&=\text{density of }N(0,R_\theta D_sR_\theta^\T),\\
 H_{\theta,s}(y)&=\prod_i\left[1+\frac{k}{A}h(z_i)\right],
 \qquad z=D_s^{-1/2}R_\theta^\T y.
 \end{aligned}
\end{equation}
Each factor of \(H_{\theta,s}\) lies in \([1/2,3/2]\).
Write \(a=0.9\), and let \(\varphi_a\) and \(\varphi_1\) be the
\(N(0,aI_d)\) and \(N(0,I_d)\) densities.
Since every eigenvalue of \(D_s\) lies in \([a,1]\), direct comparison
of Gaussian exponents and normalizing constants gives
\begin{equation}\label{eq:gaussianenvelopes}
 \begin{aligned}
 g_{\theta,s}(y)&\ge a^{d/2}\varphi_a(y),&
 g_{\theta,s}(y)&\le a^{-d/2}\varphi_1(y),\\
 p_{\theta,s}(y)&\ge c_0\varphi_a(y),&
 c_0&=2^{-d}a^{d/2}.
 \end{aligned}
\end{equation}

We next give a pointwise derivative bound, without a local-information
approximation. Let \(J\) be the skew-symmetric rotation generator,
with \(\norm J_{\op}=1\). Then \(R_\theta'=R_\theta J\), and
\[
 \partial_\theta\log g_{\theta,s}(y)
 =-\frac12y^\T R_\theta[J,D_s^{-1}]R_\theta^\T y,
\]
where \([A,B]=AB-BA\).
For \(D_s=I\), the commutator is zero. For \(D_1\) above,
\[
 \norm{D_1^{-1}-I}_{\op}
 \le\frac{dk}{1-dk}\le\frac{10d}{9}k.
\]
Since \([J,I]=0\), the absolute value of the log-density derivative is
at most \(C_dk\norm y^2\).
Also \(\norm{\partial_\theta z}\le a^{-1/2}\norm y\).
Differentiate the finite product in \eqref{eq:rotationdensity}, use
boundedness of \(h'\), and bound its remaining factors by \(3/2\).
This gives
\[
 |\partial_\theta H_{\theta,s}(y)|\le C_dk\norm y.
\]
Constants may depend on \(A\) and the fixed bounds on \(h,h'\), and are
uniform in \(k,s,\theta\).
Together with \eqref{eq:gaussianenvelopes}, these inequalities imply
\begin{equation}\label{eq:densityderivative}
 |\partial_\theta p_{\theta,s}(y)|
 \le L_dk(1+\norm y^2)\varphi_1(y).
\end{equation}

Integrating from \(\theta'\) to \(\theta\) and using the denominator
bound in \eqref{eq:gaussianenvelopes}, we obtain
\begin{align}
 \chi^2(P_{\theta,s}\Vert P_{\theta',s})
 &=\int\frac{(p_{\theta,s}-p_{\theta',s})^2}{p_{\theta',s}}\notag\\
 &\le\frac{L_d^2}{c_0}k^2|\theta-\theta'|^2
   \int(1+\norm y^2)^2\frac{\varphi_1(y)^2}{\varphi_a(y)}\,dy\notag\\
 &\le L_*^2k^2|\theta-\theta'|^2.\label{eq:chisquare}
\end{align}
The integral is finite: apart from a polynomial and a fixed normalizing
constant, its exponential factor is
\[
 \exp\left[-\left(1-\frac1{2a}\right)\norm y^2\right]
 =\exp\left[-\frac49\norm y^2\right].
\]
This also justifies the differentiations and integrations.
For \(S=0\), the Gaussian derivative vanishes and the product multiplier
alone gives the same bound. At \(k=0\), all regimes are Gaussian and
isotropic, so every rotation has the same observation law.

Let \(P_\theta^T\) be the full independent sample law with the given
regime counts. The inequality \(\log x\le x-1\) gives
\(\KL(P\Vert Q)\le\chi^2(P\Vert Q)\).
Additivity of log likelihood ratios for independent observations yields
\begin{equation}\label{eq:productKL}
 \KL(P_\theta^T\Vert P_{\theta'}^T)
 \le TL_*^2k^2|\theta-\theta'|^2.
\end{equation}
We may use \(\TV(P,Q)\le\sqrt{\KL(P\Vert Q)}\).
For a direct verification, put \(a_0=\int\sqrt{pq}\).
Jensen's inequality gives
\[
 \KL(P\Vert Q)\ge-2\log a_0
 \ge2(1-a_0)=\int(\sqrt p-\sqrt q)^2.
\]
Cauchy--Schwarz applied to
\((\sqrt p-\sqrt q)(\sqrt p+\sqrt q)\) gives
\(\TV(P,Q)^2\le2(1-a_0)\). Hence
\begin{equation}\label{eq:TV}
 \TV(P_\theta^T,P_{\theta'}^T)
 \le L_*k\sqrt T\,|\theta-\theta'|.
\end{equation}
Adding the same independent randomization device to both experiments
leaves their total variation distance unchanged. The lower bounds
therefore apply to randomized confidence procedures.

\section{Matching confidence lower bounds at every confidence level}
\label{sec:lower}
The rotation targets have a uniformly nondegenerate metric on the short
arc. For \(|\theta|,|\theta'|\le\theta_0=1/8\),
\[
 \norm{R_\theta-R_{\theta'}}_{\F}
 =2\sqrt2\left|\sin\frac{\theta-\theta'}2\right|,
 \qquad \norm{R_\theta-I}_{\F}\le\sqrt2\,\theta_0.
\]
Every nonidentity signed permutation \(P\) has
\(\norm{I-P}_{\F}\ge2\): either a diagonal sign changes, or at least
two diagonal entries vanish and their off-diagonal entries contribute
the corresponding squared norm. The triangle inequality gives
\[
 \norm{R_\theta-R_{\theta'}P}_{\F}
 \ge2-2\sqrt2\,\theta_0\quad(P\ne I),
\]
whereas the identity matching has norm at most \(2\sqrt2\,\theta_0\).
Thus the identity is the closest matching on this arc.
Using \(\sin x\ge2x/\pi\) for \(0\le x\le\pi/2\), we obtain
\begin{equation}\label{eq:rotationmetric}
 d_{\cG}([B_\theta],[B_{\theta'}])
 \ge c_{\mathrm{rot}}|\theta-\theta'|,\qquad
 c_{\mathrm{rot}}=\frac{2\sqrt2\,b_0}{\pi}>0.
\end{equation}
Here \(\pi\) denotes the circle constant.

Fix \(0<\alpha<1\), put \(p=1-\alpha\), and enlarge \(L_*\) to be
positive if necessary. Define
\[
 \eta_T=
 \begin{cases}
  \min\{\theta_0,p/(4L_*k\sqrt T)\},&k>0,\\
  \theta_0,&k=0.
 \end{cases}
\]
Every experiment with \(|\theta|\le\eta_T\) belongs to \(\cM_T(k)\),
and \eqref{eq:TV} gives \(\TV(P_\theta^T,P_0^T)\le p/4\).
At \(k=0\), this follows from exact equality of the laws.
Let \(C\) be any finite-\(T\) honest procedure in the definition of
\(\mathcal R_T(k,\alpha)\). For each such \(\theta\),
\begin{equation}\label{eq:gridcoverage}
 \Prob_0^T\{[B_\theta]\in C\}
 \ge p-\frac p4=q,\qquad q=\frac{3p}{4}.
\end{equation}
Choose an integer \(N\ge2\) with \(Nq\ge2\), and take \(N\) equally
spaced angles in \([-\eta_T,\eta_T]\), with spacing
\(\Delta_\theta=2\eta_T/(N-1)\).
Let \(Z\) count the targets contained in \(C\).
Then \(\E_0Z\ge Nq\).
If \(Z\ge1\), the largest and smallest covered indices differ by at
least \(Z-1\); \eqref{eq:rotationmetric} gives
\[
 \diam(C)\ge c_{\mathrm{rot}}\Delta_\theta(Z-1)_+.
\]
This is also true when \(C\) is empty.
Since \((Z-1)_+\ge Z-1\),
\begin{equation}\label{eq:gridlower}
 \E_0\diam(C)
 \ge c_{\mathrm{rot}}\,2\eta_T\frac{Nq-1}{N-1}
 \ge c_{\mathrm{rot}}q\eta_T.
\end{equation}
The last inequality uses \(Nq\ge2\) and \(N/(N-1)\ge1\).
Only finitely many measurable coverage events are used; no joint
measurability in a continuum parameter is presumed.

For positive \(A,B\) and \(x\ge0\),
\(\min\{A,Bx\}\ge\min\{A,B\}\min\{1,x\}\). Hence, also at \(k=0\),
\[
 \eta_T\ge
 \min\left\{\theta_0,\frac{1-\alpha}{4L_*}\right\}r(T,k).
\]
The reference model \(\theta=0\) belongs to \(\cM_T(k)\), so its
expected diameter bounds the supremum risk from below.
Taking the infimum over honest procedures proves the lower half of
Theorem~\ref{thm:main}, with
\[
 c_\alpha=c_{\mathrm{rot}}\frac{3(1-\alpha)}4
       \min\left\{\theta_0,\frac{1-\alpha}{4L_*}\right\}>0.
\]
This holds for every \(0<\alpha<1\); no restriction
\(\alpha<1/2\) was used. Combining with Section~\ref{sec:upper}
proves \eqref{eq:mainrate}.

For asymptotic honesty, set \(\alpha'=(1+\alpha)/2\).
The defining lower limit implies that every asymptotically honest
sequence eventually has coverage at least \(1-\alpha'\), uniformly
over \(\cM_T(k_T)\). Apply the finite-\(T\) argument with \(\alpha'\)
at each such \(T\). It gives \eqref{eq:asymptotic} with
\(c'_\alpha=c_{\alpha'}>0\).
The finite-sample procedures of Sections~\ref{sec:outer} and 6A are
automatically asymptotically honest. This proves the separate asymptotic
assertion with its stated quantifiers, including arbitrary weakening
rates and the identically zero sequence.\hfill\(\square\)

\section{Sharp squared-error point estimation}\label{sec:estimation}
The same geometry gives a second informativeness criterion.
We do not infer squared-error risk from coverage at a fixed confidence
level. Let \(\Theta_0\) be the compact outer set without a strength
restriction and choose a fixed countable dense sequence
\((\vartheta_j)\) in it.
The function
\[
 d_F(y)=\inf_{\vartheta\in\Theta_0}\norm{F(\vartheta)-y}
\]
is the distance to the compact set \(F(\Theta_0)\), so it is continuous
and its infimum is attained.
Choose the first index \(j\) whose residual at \(\widehat F\) is at most
\(d_F(\widehat F)+T^{-1}\).
Such an index exists by density and continuity, and the first-index
selection is Borel.
Let \(\widehat B\) be its \(B\) component.
This is an explicit measurable approximate minimum-distance rule; it does
not assume a measurable choice of an unspecified exact optimizer.

For a true tuple \(\vartheta_*\), write
\(e_*=\norm{\widehat F-F(\vartheta_*)}\).
Since \(\vartheta_*\in\Theta_0\),
\[
 \norm{F(\vartheta_j)-F(\vartheta_*)}\le2e_*+T^{-1}.
\]
Use \eqref{eq:inverseB} with the true tuple first.
For \(\delta_*\ge k>0\), \(k\le\bar k\le1\),
\[
 d_{\cG}([\widehat B],[B_*])^2
 \le\min\left\{D_{\max}^2,
       \left(\frac{L_I}{k}\right)^2(2e_*+T^{-1})^2\right\}.
\]
The inequality \((x+y)^2\le2x^2+2y^2\) and
\eqref{eq:concentration} give
\[
 \E_m d_{\cG}([\widehat B],[B_*])^2
 \le\min\left\{D_{\max}^2,
       \left(\frac{L_I}{k}\right)^2
           \left(\frac{8V}{T}+\frac2{T^2}\right)\right\}.
\]
Thus the worst-case squared-error risk is at most a fixed multiple of
\(\min\{1,(Tk^2)^{-1}\}\).
At \(k=0\), boundedness of the target gives a constant upper bound.
The estimator itself does not use \(k\).

For the matching lower bound, take the rotation models with angles
\(\pm\eta_T^{\mathrm{pt}}\), where
\[
 \eta_T^{\mathrm{pt}}=
 \begin{cases}
  \min\{\theta_0,(4L_*k\sqrt T)^{-1}\},&k>0,\\
  \theta_0,&k=0.
 \end{cases}
\]
Their total variation distance is at most \(1/2\) by \eqref{eq:TV},
and their target separation \(\Delta\) satisfies
\(\Delta\ge2c_{\mathrm{rot}}\eta_T^{\mathrm{pt}}\)
by \eqref{eq:rotationmetric}.
An estimator yields a test by assigning its estimate to the closer
target. On a test error, the estimation distance from the true target
is at least \(\Delta/2\), by the metric triangle inequality.
The sum of test-error probabilities is at least \(1-\TV\):
for the corresponding measurable event \(E\),
\[
 P_0(E)+P_1(E^c)
 =1+P_0(E)-P_1(E)\ge1-\TV(P_0,P_1).
\]
The larger squared-error risk is consequently at least
\[
 \frac{\Delta^2(1-\TV)}8\ge\frac{\Delta^2}{16}.
\]
This is a positive constant times
\(\min\{1,(Tk^2)^{-1}\}\), also at \(k=0\).
The randomization observation in Section~\ref{sec:closeness} makes the
bound valid for randomized estimators.
Hence the approximate minimum-distance rule has the sharp worst-case
squared-error rate in the same declared experiment.

\section{Explicit labels for a labelled impact matrix}\label{sec:labels}
The quotient convention describes what the statistical features identify
without additional economic labels. A labelled response needs an actual
label restriction; sorting nearly equal estimates would not give a
uniform label guarantee.
Here is a precise nonempty alternative class.
With \(b_0=(b_-+b_+)/2\), fix
\[
 0<\rho\le\min\left\{\frac{b_0}{4},\frac{b_+-b_-}{4}\right\},
\]
and restrict \(B\) to the known closed region
\[
 \mathcal K_{\mathrm{lab}}
 =\{B:\norm{B-b_0I}_{\F}\le\rho\}.
\]
This explicit loading and label restriction is imposed on both structural
models and candidates. It says that the ordered columns load near the
specified positive coordinate directions. It supplies statistical labels;
attaching economic names still requires that the restrictions have the
asserted economic interpretation. That interpretation is not inferred
solely from the observation law.

All matrices in \(\mathcal K_{\mathrm{lab}}\) have singular values between
\(b_0-\rho\) and \(b_0+\rho\), inside the original spectral interval.
For \(B,B'\) in this region, the identity matching has distance at most
\(2\rho\). For every \(P\ne I\) in \(\cG\),
\begin{align*}
 \norm{B-B'P}_{\F}
 &\ge b_0\norm{I-P}_{\F}
       -\norm{B-b_0I}_{\F}-\norm{B'-b_0I}_{\F}\\
 &\ge2b_0-2\rho>2\rho.
\end{align*}
Thus \(d_{\cG}([B],[B'])=\norm{B-B'}_{\F}\) throughout this region.
No local label-selection step remains unproved.

Intersect \(\Theta_0\) or \(\Theta_k\) with this region before constructing
confidence sets and estimators. Every concentration and inverse argument
remains valid in ordinary Euclidean matrix distance.
The lower family stays in the region after replacing \(\theta_0\) by
\[
 \min\left\{\frac18,\frac{\rho}{2b_0}\right\},
\]
since \(\norm{b_0R_\theta-b_0I}_{\F}
\le\sqrt2\,b_0|\theta|\le\rho\).
Therefore the same sharp confidence-diameter and squared-error rates
hold for the labelled impact matrix on this explicit class, at every
confidence level and for every weakening sequence.
This corollary states its label assumption; it does not assert that the
general unlabelled experiment identifies economically named shocks.

\section{Response functionals and strongly identified components}
\label{sec:responses}
The primary response in Theorem~\ref{thm:main} is the full quotient
impact matrix; Section~\ref{sec:labels} gives a labelled full-matrix
response. For a continuous response map \(\psi\) on either declared
compact target, the image of a constructed confidence set has coverage
for \(\psi(B)\). Its image is compact and is measurable as a random
closed set: its hit event for an open set is the hit event of the original
set for the open preimage.
If \(\psi\) has Lipschitz constant \(L_\psi\), then
\[
 \diam(\psi(C_T))\le L_\psi\diam(C_T).
\]
It inherits the expected-diameter upper rate and the point-estimation
upper rate, the latter with squared constant \(L_\psi^2\).
More generally, the image diameter is bounded pointwise by
\(\omega_\psi(\diam(C_T))\), where \(\omega_\psi\) is the uniform modulus
of continuity on the compact target. An upper bound alone implies no
matching lower rate.

A precise sufficient condition for the matching lower rate is a
nonflat weak-rotation response: on a short hard arc above, suppose
\[
 \norm{\psi([B_\theta])-\psi([B_{\theta'}])}
 \ge\ell_\psi|\theta-\theta'|\qquad(\ell_\psi>0)
\]
for every pair of angles in the arc.
The finite-grid proof in Section~\ref{sec:lower} applies with
\(\ell_\psi\) in place of \(c_{\mathrm{rot}}\), as does the two-point
proof in Section~\ref{sec:estimation}.
These lower rates match the preceding Lipschitz upper rates for
responses satisfying both conditions.
For labelled matrices an off-diagonal impact coordinate is an example:
\((B_\theta)_{12}=-b_0\sin\theta\), with derivative bounded away from
zero on a sufficiently short arc.
A constant response has zero estimation risk and zero confidence
diameter and does not satisfy the nonflat condition.

A distinct response is \(\Sigma_0=BB^\T\), identified by the baseline
second moment without resolving the rotation.
The confidence set
\[
 C_\Sigma=\left\{\Sigma:
     b_-^2I\le\Sigma\le b_+^2I,\quad
     \norm{\Sigma-\widehat\Sigma_0}_{\F}\le R_T\right\}
\]
has coverage at least \(1-\alpha\) by the same moment bound and has
deterministic diameter at most \(2R_T\).
Approximate minimum distance over the deterministic compact spectral set
\(\{\Sigma=\Sigma^\T:b_-^2I\le\Sigma\le b_+^2I\}\), rather than over the
possibly empty data-dependent confidence set \(C_\Sigma\), gives
squared-error risk \(O(T^{-1})\), by the residual argument of
Section~\ref{sec:estimation}.
These bounds are uniform even at \(\delta=0\).

These covariance rates are sharp in the declared classes.
For an explicit lower bound, retain the shocks of Section~\ref{sec:hard}
and replace \(B_\theta\) by
\[
 B_t=b_0e^tI,\qquad |t|\le t_0,
\]
where \(t_0>0\) is sufficiently small that the singular values stay
strictly between \(b_-\) and \(b_+\).
Shrink \(t_0\) if necessary so that \(e^{2t}\in[0.95,1.05]\).
For the labelled class, shrink \(t_0\) further until
\[
 \sqrt d\,b_0|e^t-1|\le\rho\qquad\text{throughout }|t|\le t_0.
\]
The scale family then stays in \(\mathcal K_{\mathrm{lab}}\) as well.
In coordinates divided by \(b_0\), its density is
\[
 p_t(y)=e^{-dt}p_0(e^{-t}y).
\]
Here \(p_0\) is the appropriate regime density with covariance
eigenvalues in \([0.9,1]\) and multiplier between \(2^{-d}\) and
\((3/2)^d\).
Its derivatives are bounded by a fixed polynomial times its Gaussian
envelope, uniformly in \(k\le\bar k\), because \(h,h'\) are bounded.
Differentiation gives
\[
 \partial_t p_t(y)=-dp_t(y)-y\mathbin{\cdot}\nabla_y p_t(y).
\]
These derivatives have uniformly bounded weighted \(L^2\) norm with
inverse \(N(0,0.8I)\) density: their Gaussian covariance eigenvalues lie
in \([0.855,1.05]\), the density dominates a positive multiple of
the \(N(0,0.8I)\) density, and \(1.05<2(0.8)\) ensures integrability
of the squared Gaussian envelope times any fixed polynomial.
The mean-value and denominator argument of Section~\ref{sec:closeness}
therefore proves
\[
 \chi^2(P_t\Vert P_{t'})\le C|t-t'|^2
\]
in every regime. The \(T\)-sample total variation is at most
\(C\sqrt T\,|t-t'|\).
The covariance curve \(b_0^2e^{2t}I\) has Frobenius distance bounded
below by a positive constant times \(|t-t'|\) on this interval.
Apply the finite-grid proof on an interval of radius of order
\(T^{-1/2}\), and the two-point proof to two such parameters.
They give, respectively, expected diameter at least
\(c_\alpha T^{-1/2}\) and squared-error risk at least \(c/T\).
The nuisance features remain those of Section~\ref{sec:hard} and meet
\(\delta\ge k\). These are lower bounds inside \(\cM_T(k)\).

Every Lipschitz response depending only on \(\Sigma_0\) inherits the
strong upper rate. A matching lower rate for a particular response
requires a corresponding nonflat covariance curve.
Rotation-sensitive, covariance-only, and flat responses must therefore
be distinguished; the theorem does not assign the weak-orientation rate
to every possible \(\psi\).

\section{Dimension one and two checks of the joint index}\label{sec:checks}
When \(d=1\), there is no rotational identification problem.
Choose the column sign explicitly by \(B>0\).
In the baseline regime, \(\E E_0^2=1\), so \(\E U^2=B^2\).
The baseline second-moment confidence interval and
\[
 |B-B'|\le\frac{|B^2-(B')^2|}{2b_-}
\]
give finite-sample uniform coverage and diameter \(O(T^{-1/2})\)
for \(B\) throughout the same bounded-moment class.
Approximate minimum distance gives squared-error risk \(O(T^{-1})\).
For the lower bound, choose Gaussian shocks in every regime, \(D_s=1\),
and \(B=b_0e^t\) in an interior interval.
Their eighth moments equal \(105\), so they are permitted.
The scale-density calculation in Section~\ref{sec:responses}, together
with the finite-grid and two-point arguments, gives matching lower
bounds for all confidence levels.
The pairwise index \(\delta\) is neither defined nor needed in dimension
one; no one-shock weak-rotation rate is asserted.

The joint index can be positive even when neither feature collection
identifies on its own. For \(d=3\), \(S\ge1\), and small \(k>0\), take
\[
 c=(0,k,0),\qquad D_1=\diag(1,1,1-k),
\]
and \(D_s=I\) for all other regimes.
Use density \(f_{k/A}\) for shock \(2\), Gaussian densities for shocks
\(1,3\), and scale the third shock by \(\sqrt{1-k}\) in regime \(1\).
Every coordinate eighth moment is at most \(105\).
The variance profiles of shocks \(1,2\) coincide, while the fourth
cumulants of shocks \(1,3\) are both zero.
Each of the separate moment maps thus leaves a rotation unresolved.
The joint squared pair margins are nevertheless \(k^2,k^2,2k^2\), so
\(\delta=k\), and the global inverse bound applies.
This example illustrates the arbitrary-\(d\) theorem and does not
replace it.

Finally, zero fourth cumulant does not imply Gaussianity.
A variable taking values \(\sqrt3\) and \(-\sqrt3\), each with
probability \(1/6\), and \(0\) with probability \(2/3\), has mean zero,
variance \(1\), fourth moment \(3\), sixth moment \(9\), and eighth
moment \(27\). It is non-Gaussian although its fourth cumulant is zero.
Such laws are allowed by the model class.
The constant worst-case lower bound at \(k=0\) uses the rotationally
indistinguishable Gaussian subfamily. It does not assert that every
individual \(\delta=0\) law is unidentified from its full distribution,
or that higher-order features could not improve inference at that law.

\section{Scope, verification, and the canonical question}\label{sec:scope}
The canonical C32-80 question requires declarations of the structural
class \(\cM_T\), observation and dependence model, moment assumptions,
signs and labels, response, and informativeness criterion.
Sections~\ref{sec:model} and \ref{sec:labels} supply these declarations.
The core experiment allows arbitrary unknown conditional shock
distributions under coordinate independence and an eighth-moment bound.
Normal or parametric laws are used only as explicit lower-bound
subfamilies. The full impact response with its sign and permutation
convention is the primary response; the labelled and other-response
results have their conditions stated separately.

Theorem~\ref{thm:main} provides finite-sample uniform coverage, sharp
worst-case expected diameter as the variance and fourth-feature lower
bound shrinks, and the corresponding conclusion for the exact
asymptotic coverage quantifier in the canonical statement.
Section~6A provides a single procedure that does not need to know that
lower bound. Section~\ref{sec:estimation} gives matching minimax
squared-error bounds. The nonshrinking case and every sequence \(k_T\)
are included. Observationally equivalent targets, moment weakness,
unknown shock laws, and label ambiguity are not replaced by a
point-identified Gaussian model.

Known regime dates, observed innovation vectors, a common \(B\),
independence across observations, the stated eighth-moment bound, and
fixed \(d\) are explicit parts of the declared experiment.
Unknown regimes, changing \(B\), estimated VAR residuals, heavy-tail
classes without the required moments, and other non-Gaussian-feature
indices are different models; no theorem for them is claimed without
their own assumptions. A universal response-independent rate would
also be false, as the constant and covariance responses demonstrate.
The results address the fully declared structural-inference problem
permitted by the canonical agenda, rather than a classification of
every model or every source of non-Gaussian information in the literature.

Lewis~\cite{Lewis2025}, especially Sections~5 and 7.2, supplies the
identification-robust agenda and the normalization and label warnings.
The experiment and the proofs of the global inverse bound and uniform
minimax theorem are given here; that source is not represented as
already proving this exact theorem.

All arguments are hand derivations. No finite simulation, numerical
experiment, Lean formalization, external peer review, or novelty claim
is used in their place. The complete mathematical argument received
independent internal source and scope review. The accompanying review
records identify the checked versions and the scope of those checks.

\begin{thebibliography}{1}
\bibitem{Lewis2025}
Daniel J. Lewis.
\newblock Identification Based on Higher Moments in Macroeconometrics.
\newblock \emph{Annual Review of Economics} \textbf{17} (2025), 665--693.
\newblock \href{https://doi.org/10.1146/annurev-economics-070124-051419}
{doi:10.1146/annurev-economics-070124-051419}.
\newblock \href{https://discovery.ucl.ac.uk/id/eprint/10208055/1/Lewis_annurev-economics-070124-051419.pdf}
{Published article in the author's institutional repository}.
\newblock Agenda locators: Sections~5 and 7.2.
\end{thebibliography}
\end{document}
